\section{Timing, resources and usability: methods before results}
No reference-machine end-to-end latency distribution or tracing-overhead
measurement has run. Existing compiler/service timings omit important parts
of interaction. Debug native-harness execution is not a substitute for a
release application's input-to-visible-result measurement.

\subsection{Implemented optional instrumentation}
Ordinary builds disable diagnostics. A separate build with
\code{VITE\_TEXLOOM\_PERFORMANCE=1} records static categories, sequence numbers
and monotonic \code{performance.now()} offsets. The clock's purpose and its
resolution limits follow the High Resolution Time API~\cite{hrtime}.
The local collector records no source text, project paths, selected content
or error strings and transmits no trace. It exposes snapshot/stop/clear only
in the opted-in build.

Table~\ref{tab:timing} states the actual boundaries. A method-entry timestamp
is not physical input arrival. Source feedback starts after Monaco has already
changed its model. A compile trace begins only when the active compile call
starts; its preceding automatic debounce or queue delay is outside that trace.

\begin{table}[htbp]
\small\begin{tabularx}{\linewidth}{@{}>{\raggedright\arraybackslash}p{.24\linewidth}YY@{}}
\toprule
Category & Start and operation boundary & Important exclusion\\
\midrule
Source feedback & Monaco change callback to feedback/source-change update & OS dispatch and prior Monaco work\\
Source selection & Eligible selection call to resolved proof and frame & Cache state must be distinguished; not a hardware click timestamp\\
Existing-object selection & Eligible object call to frame/focus & Does not include creating its original source wrapper\\
Drag preview & First meaningful pointer-move callback to outline update & One sample per drag; no sustained frame-rate result\\
Text move/edit & Release handler or method entry to commit/publication/reselection & Physical display and prior event dispatch remain outside it\\
Manual/auto compile & Active call before source synchronization to PDF publication & Prior debounce/queue; lazy publication alone does not prove visible-page paint\\
\bottomrule
\end{tabularx}
\caption{Implemented instrumentation boundaries, not measured performance.}
\label{tab:timing}
\end{table}

Milestones include source preparation, native build return, PDF.js preparation,
frontend geometry validation, receipt recheck, native commit and publication.
The native build return already includes independent layout approval for
planned visual operations. It cannot isolate TeX from worker/transport cost.
The later receipt recheck is not another independent PDF comparison.

After publication the collector can wait for two animation callbacks while
the document is visible and identity remains current. Animation callbacks
are defined by the HTML rendering API~\cite{animationframes}; two callbacks
do not establish hardware scanout, absence of occlusion or human perception.
The field is named a \emph{frame proxy}. Long-document lazy publication receives
no such result when visible-page completion has not been established.
Operation duration and frame-proxy duration remain separate.

\subsection{Incomplete and lost traces}
The collector holds at most 256 completed records, 16 active operations and
64 milestones per operation. Evictions, dropped starts and truncated phases
are counted. Export bounded batches and preserve their fingerprints before
resetting. Explicit clear discards diagnostics; it is not a way to improve an
evaluation denominator.

Hidden, superseded, timed-out, unavailable, invalid-clock or stopped frame
observations have no completed frame-proxy latency. The presentation timeout
is 2 s and does not cancel the compilation. Technical outcomes such as
published, refused, stale, cancelled and compile-failed are not independent
correctness labels. A stopped unfinished operation remains incomplete rather
than becoming a pass.

The next, unbuilt source revision exports schema 2 with a local random
per-launch collector ID, clear epochs and issued-ID boundaries. Explicit
discards, invalid-clock starts and cumulative cleared-row counts are retained.
A missing random-ID facility leaves identity unresolved; no project identifier
is substituted. Ordinary disabled collection does not request randomness.
These changes have not been tested, and old schema-1 traces are not relabelled
as new exports.

An unexecuted analysis draft reads hashed snapshot files under an explicit
controlled root. Planned trial IDs, warmup roles and missing traces are
separate from collector record IDs. Overlapping snapshots cannot count the
same retained record twice. Context groups separate source/executable,
protocol, environment, workload, document, engine, language, cache state and
operation category. Loss, pending work, unclaimed main-category records or
missing planned samples suppress distributions for the affected current
epochs. Clear history does not prove earlier epochs were exported.

For complete current epochs the proposed summaries separate outcomes and
endpoints. Median and nearest-rank p95 describe only observed samples;
censored frame states retain counts, never zero-filled latency. A conditional
frame p95 does not describe every requested action. The analysis computes
neither independent correctness nor physical presentation or tracing overhead.
Its source and specification are in \path{research/tools/timing_records.py}
and \path{research/TIMING-RECORDS.md}; no invocation has run.

Previously recorded headless checks cover clock sampling disabled in ordinary mode, bounded storage,
failure/pending accounting and real-shell milestone plumbing with synthetic
animation callbacks. Review found and corrected an early publication freshness
check that ran before the busy flag was released. None of those diagnostics
measures an actual Windows presentation boundary or acceptable overhead.
Those earlier checks do not validate the source-only schema-2 changes.

\subsection{Required timing and resource experiment}
First validate boundaries against actual actions on the exact ordinary and
diagnostic executables. Use an external measurement method for their paired
overhead comparison; the uninstrumented build cannot supply the same internal
trace. Record release optimizations, source hashes, WebView2, CPU/RAM, fonts,
engine, DPI/display/power configuration and background load.

The proposed workload uses at least ten representative small documents, ten
declared warmups and 30 measured warm repetitions per document/engine condition,
with randomized paired order. Keep cold first-selection cases, cache state,
failed/censored attempts, sample counts, medians and p95. First-wrap movement
can include private identity work, so count logical operations, private builds
and TeX passes separately. Report source-selection cache hits and prepared
object selections as different workloads rather than infer support from speed.

A controlled single-instance long-session workload is also needed. The older
20-minute process-tree sample included shared WebView descendants from two
instances and another crash/relaunch; it is not an isolated leak or resource
stability pass. New application-opening/resource tests are deferred while the
owner is away.

\subsection{Usability is a separate question}
An unassisted task should ask a new user to identify a target, understand an
expanded selection, move/edit it, respond to a refusal, undo, save and reopen a
synthetic project. Record task success, errors, assistance and the user's
understanding of preview versus verified state. Scripted interaction does not
establish these outcomes. No participant study or unassisted usability result
is claimed in this report.
